\documentclass[10pt,a4paper]{amsart}
\usepackage[margin=19mm]{geometry}
\usepackage{amsmath,amssymb,mathtools}
\usepackage[scr=rsfs,cal=euler]{mathalfa}
\usepackage{xfrac}
\usepackage[unicode,hidelinks,bookmarks=false]{hyperref}
\newcommand{\ie}{i.e.\ }
\newcommand{\cf}{cf.\ }
\newcommand{\resp}{respectively\ }
\newcommand{\Subsection}{\subsection}
\providecommand{\nobreakdash}{\nobreak\hbox{-}}
\setlength{\emergencystretch}{2em}
\multlinegap0pt
\usepackage{bm}
\usepackage{bbm}
\usepackage{dsfont}
\usepackage{xspace}
\usepackage{cancel}
\usepackage{mathtools}
\usepackage{amsthm}
\makeatletter
\@ifundefined{theorem}{\newtheorem{theorem}{Theorem}}{}
\makeatother
\title[Cross representations of additive complements of $r$-th powe]{Cross representations of additive complements of $r$-th powers}
\author{Yuchen Ding, Ben Krause, Csaba S\'andor, Yu-Chen Sun, Zihan Zhang}
\date{Source: arXiv:2512.15407v6 (CC BY 4.0)}
\begin{document}
\maketitle

\noindent\emph{Source and licence.} Cross representations of additive complements of $r$-th powers. Version arXiv:2512.15407v6 (CC BY 4.0), distributed under the Creative Commons Attribution 4.0 International licence, as recorded in the arXiv licence metadata for this exact version and shown on the abstract page https://arxiv.org/abs/2512.15407v6. Full rights notice: licenses/arxiv-2512.15407-CC-BY-4.0.txt.\par\smallskip

\noindent\textit{Editorial scope.} Counted unit: the Theorem whose source label is \texttt{thm:3} in the article (source0.tex lines 799--804, source label \texttt{thm:3}), delivered together with the article's own complete original proof of it (source0.tex lines 805--906, one \texttt{proof} environment of 102 lines). No mathematics is written, repaired or re-derived: statement and proof are the article's own text line for line. Required context retained uncounted: eq-1216-1 (lines 139--141). Every definition, display and label of those lines is retained unchanged. Omitted from this package and not redistributed: 1--138, 142--798, 907--1001. Nothing inside a retained excerpt is omitted or altered: every retained body is delivered exactly as the article's lines are written, so no omission receipt of any kind accompanies this package. Citations: the retained excerpts carry no citation command at all, so no bibliography entry of the article is transcribed and no reference is redistributed. Reference adaptation: none was needed and none was made; every reference inside the delivered unit resolves inside it, so no unresolved reference or citation is produced. Printed slips kept literal: no printed word, symbol, hypothesis or number of the counted statement or of its proof was altered. Layout and class: the article's own class is \texttt{amsart} with packages including amsmath, amsfonts, amssymb, amsthm, latexsym, enumerate, url, cases, mathrsfs, hyperref, etoolbox; the wrapper is the lane's pinned \texttt{amsart} editorial wrapper at 10pt on A4, which restates the theorem-like environments the retained text needs and adds the article's own macro definitions that the retained text needs (none), with extra wrapper packages (bm, bbm, dsfont, xspace, cancel, mathtools). The delivered numbering is the unit's own. Assets: no retained span contains a figure environment, a graphics-inclusion command, a PDF-page inclusion, a table or a TikZ picture; no image file is copied into this package, the package declares no asset dependency and no image file is redistributed with it. Licence: version arXiv:2512.15407v6 of this article is distributed under the Creative Commons Attribution 4.0 International licence, as recorded in the arXiv licence metadata for this exact version and shown on the abstract page https://arxiv.org/abs/2512.15407v6. A full rights notice is delivered at licenses/arxiv-2512.15407-CC-BY-4.0.txt. Characters: every retained excerpt is pure ASCII, so the delivered engine, XeLaTeX with its default Latin Modern text font, reports no missing character.
\begin{align}\label{eq-1216-1}
\liminf_{N\rightarrow\infty}\mathcal{W}\big(N\big)/\sqrt{N}\ge \frac{4}{\pi}, 
\end{align}

\begin{theorem}\label{thm:3}
If $\mathcal{W}=\{w_n\}_{n=1}^{\infty}$ is an additive complement of the squares that is exact on average, then 
$$
w_n \sim \frac{\pi^2 n^2}{16},  \quad \text{or~equivalently} \quad \mathcal{W}(N)\sim \frac{4}{\pi}\sqrt{N}.
$$
\end{theorem}

\begin{proof}
Throughout the proof, let $N$ be sufficiently large. By (\ref{eq-1216-1}) we have
\begin{align*}
\limsup_{N\rightarrow\infty}\mathcal{W}\big(N\big)/\sqrt{N}\ge \liminf_{N\rightarrow\infty}\mathcal{W}\big(N\big)/\sqrt{N}\ge \frac{4}{\pi}.
\end{align*}
It remains to rule out the following two cases.

{\it Case I.} 
$$
\limsup_{N\rightarrow\infty}\frac{\mathcal{W}(N)}{\sqrt{N}}>\sqrt{6}.
$$
In this case, there are infinitely many positive integers $N$ such that
$
\mathcal{W}(N)>\sqrt{6N}.
$
For these $N$, we clearly have
$$
\sum_{n\le 2N}f(n)\ge \sum_{w\le N,\ m^2\le N}1\ge \sqrt{6N}\cdot \big(1+o(1)\big)\sqrt{N}>\sqrt{5}N,
$$
from which it follows immediately that
$$
\sum_{n\le 2N}f(n)-2N\ge (\sqrt{5}-2)N.
$$
This contradicts the assumption that $\mathcal{W}$ is exact on average.

{\it Case II.} 
\begin{align*}
\frac{4}{\pi}<\limsup_{N\rightarrow\infty}\frac{\mathcal{W}(N)}{\sqrt{N}}=\gamma\le \sqrt{6}.
\end{align*}
In this case, set 
$$
\delta=\frac{\gamma+4/\pi}{2}.
$$
Then there are infinitely many positive integers $N$ such that
\begin{align}\label{eq-1215-1}
\mathcal{W}(N)>\delta \sqrt{N}.
\end{align} 
For these $N$, let 
$$
\delta_1=\frac{\delta+4/\pi}{2} \quad \text{and} \quad \varepsilon=\left(\frac{\delta-\delta_1}{4}\right)^2.
$$ 
We further split {\it Case II} into two subcases.

{\it Subcase II-1.}  
\begin{align}\label{eq-1215-2}
\mathcal{W}\Big((1-\varepsilon)N\Big)\le\delta_1 \sqrt{N}.
\end{align} 
In this subcase we have 
\begin{align}\label{eq-1220-1}
\mathcal{W}(N)-\mathcal{W}\Big((1-\varepsilon)N\Big)>\big(\delta-\delta_1\big)\sqrt{N}
\end{align} 
by (\ref{eq-1215-1}) and (\ref{eq-1215-2}). Hence, from (\ref{eq-1220-1}) we get
\begin{align*}
\sum_{(1-\varepsilon)N<n\le (1+\varepsilon)N}f(n)&\ge \sum_{\substack{(1-\varepsilon)N<w\le N\\ m^2\le \varepsilon N}}1\\
& \ge \big(\delta-\delta_1\big)\sqrt{N}\cdot\Big(1+o(1)\Big)\sqrt{\varepsilon N}\\
&>\frac{(\delta-\delta_1)^2}{5}N.
\end{align*}
Therefore, it follows that
\begin{align*}
\sum_{n\le (1+\varepsilon)N}f(n)-(1+\varepsilon)N&\ge \sum_{(1-\varepsilon)N<n\le (1+\varepsilon)N}f(n)-2\varepsilon N+O(1)\\
&>\frac{(\delta-\delta_1)^2}{5}N- 2\left(\frac{\delta-\delta_1}{4}\right)^2N+O(1)\\
&>\frac{(\delta-\delta_1)^2}{20}N.
\end{align*}
Hence, in {\it Subcase II-1}, $\mathcal{W}$ cannot be exact on average. 

{\it Subcase II-2.}  
\begin{align}\label{eq-1215-3}
\mathcal{W}\Big((1-\varepsilon)N\Big)>\delta_1 \sqrt{N}.
\end{align} 
In this subcase, we first note that
\begin{align*}
\sum_{n\le N}f(n)=\sum_{w+m^2\le N}1=\sum_{m<\sqrt{N}}\sum_{w\le N-m^2}1=\sum_{m<\sqrt{N}}\mathcal{W}\big(N-m^2\big).
\end{align*} 
We now split the sum according as $m\le \sqrt{\varepsilon N}$ or $m>\sqrt{\varepsilon N}$. Then
\begin{align}\label{eq-1215-4}
\sum_{n\le N}f(n)=\sum_{m\le \sqrt{\varepsilon N}}\mathcal{W}\big(N-m^2\big)+\sum_{\sqrt{\varepsilon N}< m<\sqrt{N}}\mathcal{W}\big(N-m^2\big).
\end{align} 
For $m\le \sqrt{\varepsilon N}$ we have $N-m^2\ge (1-\varepsilon)N$, and hence by (\ref{eq-1215-3}) we get
\begin{align}\label{eq-1215-5}
\mathcal{W}\big(N-m^2\big)>\delta_1 \sqrt{N}.
\end{align}
By (\ref{eq-1216-1}), (\ref{eq-1215-4}) and (\ref{eq-1215-5}) we obtain 
\begin{align}\label{eq-1215-6}
\sum_{n\le N}f(n)&>\delta_1 \sqrt{N}\sum_{m\le \sqrt{\varepsilon N}}1+\sum_{\sqrt{\varepsilon N}< m<\sqrt{N}}\mathcal{W}\big(N-m^2\big)\nonumber\\
&\ge \big(\delta_1-4/\pi\big) \sqrt{\varepsilon}N +\sum_{m<\sqrt{N}}\Big(4/\pi+o(1)\Big)\sqrt{N-m^2}+O\big(\sqrt{N}\big).
\end{align}
A routine partial summation gives
\begin{align}\label{eq-1215-7}
\sum_{m<\sqrt{N}}\Big(4/\pi+o(1)\Big)\sqrt{N-m^2}\sim N \quad \text{as }N\rightarrow\infty.
\end{align}
Inserting (\ref{eq-1215-7}) into (\ref{eq-1215-6}), we then obtain
\begin{align*}
\sum_{n\le N}f(n)-N&\ge \Big(\delta_1-4/\pi+o(1)\Big) \sqrt{\varepsilon}N\\
&> \frac{\big(\delta-\delta_1\big)\big(\delta_1-4/\pi\big)}{5}N.
\end{align*} 
Hence, in {\it Subcase II-2}, $\mathcal{W}$ cannot be exact on average.
Thus,
\[
 \limsup_{N\to\infty}\frac{\mathcal{W}(N)}{\sqrt N}\le \frac{4}{\pi}.
\]
Together with (\ref{eq-1216-1}), this gives $\mathcal{W}(N)\sim \frac{4}{\pi}\sqrt N$, which is equivalent to $w_n\sim \frac{\pi^2}{16}n^2$.
\end{proof}

\end{document}
